%!TEX TS-program = pdflatex
%!TEX encoding = UTF-8 Unicode
% Software Engineer Resume -- Ssohrab Borhanian
% Requires only standard LaTeX packages. Compile with pdflatex or xelatex.

\documentclass[10pt, a4paper]{article}

% ── Encoding & fonts ──────────────────────────────────────────────────────────
\usepackage[T1]{fontenc}
\usepackage[utf8]{inputenc}
\usepackage{lmodern}          % improved vector version of Computer Modern
\usepackage{fontawesome5}

% ── Page layout ───────────────────────────────────────────────────────────────
\usepackage[a4paper,
            left=1.8cm, right=1.8cm,
            top=1.4cm,  bottom=1.6cm]{geometry}

% ── Typography ────────────────────────────────────────────────────────────────
\usepackage{microtype}

% ── Colors ────────────────────────────────────────────────────────────────────
\usepackage[dvipsnames]{xcolor}
\definecolor{accent}{RGB}{25, 95, 170}      % deep blue for headings/links
% \definecolor{accent}{RGB}{0, 0, 0}      % deep blue for headings/links
\definecolor{muted}{RGB}{90, 90, 90}        % dark gray for secondary text

% ── Hyperlinks ────────────────────────────────────────────────────────────────
\usepackage[colorlinks=true,
            urlcolor=accent,
            linkcolor=accent]{hyperref}

% ── Section headings ──────────────────────────────────────────────────────────
\usepackage{titlesec}
\titleformat{\section}
  {\color{accent}\large\bfseries}
  {}{0em}{\MakeUppercase}
  [\vspace{-8pt}\color{accent}\rule{\linewidth}{0.5pt}]
\titlespacing*{\section}{0pt}{7pt}{3pt}

% ── Lists ─────────────────────────────────────────────────────────────────────
\usepackage{enumitem}
\setlist[itemize]{
  leftmargin=1.6em,
  itemsep=0pt,
  parsep=0pt,
  topsep=2pt,
  label=\small$\bullet$
}

% ── Tables ────────────────────────────────────────────────────────────────────
\usepackage{tabularx}
\usepackage{array}

% ── Misc ──────────────────────────────────────────────────────────────────────
\setlength{\parindent}{0pt}
\setlength{\parskip}{0pt}
\pagestyle{plain}

%──────────────────────────────────────────────────────────────────────────────
% CUSTOM COMMANDS
%──────────────────────────────────────────────────────────────────────────────

% \cventry{Role}{Organization}{Location}{Date range}
\newcommand{\cventry}[4]{%
  \textbf{#1}                      \hfill {\color{muted}\small\textit{#4}} \\[2pt]
  {\color{muted}\small\textit{#2}} \hfill {\color{muted}\small\textit{#3}}
  % \begin{tabularx}{\linewidth}{@{}X @{}r}
  %   \textbf{#1}                      & {\color{muted}\small\textit{#4}} \\
  %   {\color{muted}\small\textit{#2}} & {\color{muted}\small\textit{#3}}
  % \end{tabularx}%
}

% \cvskill{Category}{Values}
\newcommand{\cvskill}[2]{%
  \noindent
  % \begin{tabularx}{\linewidth}{@{}>{\bfseries\small}p{3.4cm} X@{}}
    % #1 & {\small #2}
    \begin{tabularx}{\linewidth}{@{}p{3.2cm} X@{}}
    \textbf{\small #1} & {\small #2}
  \end{tabularx}\raggedright%
}

\newcommand{\tpc}{\textperiodcentered}

\pagenumbering{gobble}

%──────────────────────────────────────────────────────────────────────────────
\begin{document}
%──────────────────────────────────────────────────────────────────────────────

% ── HEADER ───────────────────────────────────────────────────────────────────
\begin{center}
  {\Huge\bfseries Ssohrab Borhanian, Ph.D.}\\[3pt]
  % {\large\color{accent} Scientific Software Engineer \,$\cdot$\, Python Developer}\\[3pt]
  {\small\color{muted}
    % Steglitzer Damm 113E, 12169 Berlin, Germany
    % ~\tpc~
    % \href{tel:+4915251705163}{+49\,152\,51705163}
    % ~\tpc~
    \faEnvelope~\href{mailto:sborhanian@gmail.com}{sborhanian@gmail.com}
    ~|~
    \faGlobe~\href{https://sborhanian.github.io/}{sborhanian.github.io}
    ~|~
    \href{https://www.linkedin.com/in/ssohrab-borhanian}{\faLinkedin}
    ~
    \href{https://github.com/sborhanian} {\faGithub}
    ~
    \href{https://gitlab.com/sborhanian} {\faGitlab}
    ~|~
    EU Citizen
  }
\end{center}

\vspace{1pt}

% ── PROFILE ──────────────────────────────────────────────────────────────────
\section{Executive Summary}
\begin{itemize}
  \item Open-source, scientific software developer with background in Python and ML tools.
  \item Physicist with 10+ years of academic expertise in gravitational-wave astrophysics and large-scale collaborations.
  \item Avid collaborator and communicator with strong academic publication and presentation record.
  \item Problem-solver and keen learner at heart; highly adaptable to new environments and challenges.
\end{itemize}

% ── EDUCATION ────────────────────────────────────────────────────────────────
\section{Education}

\cventry%
  {The Pennsylvania State University }%
  {Ph.D.\ and M.Sc.\ in Physics~\tpc~ Fulbright Scholar~\tpc~ Peter Clay Eklund Lecturing Award}%
  {State College, PA, USA}%
  {2014 -- 2021}

  % \begin{itemize}
  %   \item {\small Dissertation: \textit{Compact Binaries as Astrophysical Laboratories and the Promise of Third-Generation Detectors}}
  % \end{itemize}

\vspace{6pt}
\cventry%
  {Technical University of Berlin}%
  {B.Sc.\ in Physics and Mathematics ~\tpc~ Only double degree in my cohort}%
  {Berlin, Germany}%
  {2009 -- 2014}
  
  % \begin{itemize}
  %   \item {\small Physics specialization: Astrophysics --- Mathematics specialization: Differential Geometry}
  % \end{itemize}

% ── OPEN SOURCE PROJECTS ─────────────────────────────────────────────────────
\section{Open Source Projects}

\cventry%
  {\textsc{gwbench} --- Gravitational-Wave Detector and Network Benchmarking}%
  {Lead Developer and Maintainer ~\tpc~ 
   \href{https://gitlab.com/sborhanian/gwbench}{\faGitlab}  ~\tpc~ 
   \href{https://arxiv.org/abs/2010.15202}{\textit{Class.~Quant.~Grav.}~38,~2021, [arXiv:2010.15202]} }%
  {}%
  {2019 -- Present}
\begin{itemize}
  \item Open-source, Python package for Fisher-Matrix-based benchmarking of gravitational-wave detector networks.
  \item Adopted as a primary analysis tool by the Einstein Telescope and Cosmic Explorer collaborations.
  \item Highly modular and streamlined user experience for large-scale computational studies.
\end{itemize}

\vspace{3pt}
\cventry%
  {\textsc{xkn} --- Kilonova Modeling Framework}%
  {Core Developer ~\tpc~ \href{https://github.com/GiacomoRicigliano/xkn}{\faGithub}
   ~\tpc~
   \href{https://arxiv.org/abs/2311.15709}{\textit{Mon.~Not.~Roy.~Astron.~Soc.}~529~(2024)~1,~647-663,, [arXiv:2311.15709]} }%
  {}%
  {2021 -- 2023}
\begin{itemize}
  \item Python package to generate kilonova light curves from binary neutron star mergers.
  \item {\em \color{muted} Key contributions}: Implementation of accessible user interface and optimization for computational efficiency.
\end{itemize}

% ── EXPERIENCE ───────────────────────────────────────────────────────────────
\section{Experience}

\cventry%
  {Postdoctoral Researcher}%
  {Universit\`a degli Studi di Milano-Bicocca
  ~\tpc~ Research group: Prof.\ Davide\ Gerosa}%
  {Milan, Italy}%
  {09/2023 -- Present}
\begin{itemize}
  \item Developing \textsc{jax}-based ML-pipeline for gravitational-wave data analysis.
  \item Extending the Python package \href{https://github.com/dgerosa/precession}{\textsc{precession}} to support \texttt{jit}-compilation and GPU-acceleration through \textsc{jax}.
  \item Investigated the impact of facility timing on the science case of third-generation, gravitational-wave detectors.
  % \item Developing a ML-based data analysis pipelines for long-duration gravitational-wave signal detection and parameter estimation
  % \item Developing a fast and efficient, non-parametric framework for astrophysical population inference of gravitational-wave sources
  % \item Assessed of impact of facility timing uncertainties on the science case of third-generation gravitational-wave detectors
\end{itemize}

\vspace{6pt}
\cventry%
  {Postdoctoral Researcher}%
  {Friedrich-Schiller-Universit\"at Jena ~\tpc~ Research group: Prof.\ Sebastiano\ Bernuzzi}%
  {Jena, Germany}%
  {06/2021 -- 08/2023}
\begin{itemize}
  \item Contributed to and extended the functionality of the open-source Python packages \textsc{gwbench}, \textsc{xkn}, and \href{https://github.com/matteobreschi/bajes}{\textsc{bajes}}.
  % \item Extended \textsc{gwbench}: numerical improvements, bug fixes, and expansion of the functionality via additional features
  % \item Contributed as a core developer to the \textsc{xkn} Python package for kilonova light-curve modeling from binary neutron star mergers
  % \item Contributed to core modules of the \textsc{bajes} Python code for Bayesian inference of gravitational-wave sources
\end{itemize}

\vspace{6pt}
\cventry%
  {Graduate Researcher (Ph.D.)}%
  {The Pennsylvania State University ~\tpc~ Advisor: Prof.\ Bangalore S.\ Sathyaprakash}%
  {State College, PA, USA}%
  {08/2016 -- 05/2021}
\begin{itemize}
  \item Developed \textsc{gwbench} as a modular Python framework for large-scale science case studies.
  \item Investigated the science case of third-generation, gravitational-wave detectors via comprehensive trade studies.
  \item Investigated numerical relativity waveforms for a post-Newtonian signature in the strong-field regime.
  % \item Designed and built \textsc{gwbench} from scratch as a modular Python framework
  %       for Fisher-information-based gravitational-wave detector benchmarking and science case studies
  % \item Conducted a comprehensive computational trade study of third-generation
  %       gravitational-wave detector geometries and network configurations for the Einstein Telescope and Cosmic Explorer
  % \item Developed numerical and analytic models for binary black hole waveforms and
  %       tests of general relativity
  % \item Member of the LIGO Scientific Collaboration and Cosmic Explorer Consortium
\end{itemize}

\vspace{6pt}
\cventry%
  {Research Intern}%
  {Duke University~\tpc~ RISE Worldwide Scholarship~\tpc~ Advisor: Prof.\ Kate\ Scholberg}%
  {Durham, NC, USA}%
  {08/2012 -- 09/2012}
% \begin{itemize}
%   \item Implemented numerical simulations of pre-supernova neutrino emission as an
%         early-warning signal for stellar collapse events
% \end{itemize}

% ── SKILLS ───────────────────────────────────────────────────────────────────
\section{Technical Skills}
\cvskill{Programming}{Python, Swift, SwiftUI, Fortran, Java, Mathematica, MATLAB.}
\cvskill{Python Ecosystem}{astropy, h5py, jax, joblib, matplotlib, mpmath, numpy, pandas, scipy, sympy.}
\cvskill{Developer Tools}{Git/GitHub/GitLab, HPC clusters (SLURM/HTCondor), \LaTeX{}.}
\cvskill{Data Analysis}{Machine Learning, Bayesian inference, Monte Carlo methods, time series analysis.}

\section{Soft Skills}
\cvskill{Communication}{20+ presentations and invited seminars at international conferences and universities.}
\cvskill{Technical Writing}{\href{https://inspirehep.net/authors/1730808}{25 academic publications} (3700+ citations).}
\cvskill{Languages}{\textit{Fluent:} German, English, Persian; \quad \textit{Basic:} French, Italian.}


% ── SELECTED HONORS ──────────────────────────────────────────────────────────
% \section{Selected Honors}

% \begin{tabularx}{\linewidth}{@{}X r@{}}
%   \textbf{Peter Clay Eklund Lecturing Award},
%     \textit{\small Department of Physics, Pennsylvania State University}
%     & \small 2021 \\[3pt]
%   \textbf{Fulbright Scholarship},
%     \textit{\small German-American Fulbright Commission}
%     & \small 2014 \\[3pt]
%   \textbf{RISE Worldwide Scholarship},
%     \textit{\small German Academic Exchange Service (DAAD)}
%     & \small 2012 \\[3pt]
%   \textbf{Acknowledgement for Excellent Achievements in Physics},
%     \textit{\small German Physical Society}
%     & \small 2009 \\
% \end{tabularx}

% ── SELECTED PUBLICATIONS ────────────────────────────────────────────────────
% \section{Selected Publications}

% \noindent\textbf{Borhanian} ---
% \textsl{\textsc{gwbench}: a novel Fisher information package for gravitational wave benchmarking}\\
% {\small \textit{Class.~Quant.~Grav.}~38~(2021)~17,~175014
% \hfill \href{https://arxiv.org/abs/2010.15202}{arXiv:2010.15202}}

% \vspace{3pt}
% \noindent\textbf{Borhanian}, Sathyaprakash ---
% \textit{Listening to the Universe with Next Generation Ground-Based GW Detectors}\\
% {\small \textit{Phys.~Rev.~D}~110~(2024)~8,~083040
% \hfill \href{https://arxiv.org/abs/2202.11048}{arXiv:2202.11048}}

% \vspace{3pt}
% \noindent\textbf{Borhanian}, Dhani, Gupta, Arun, Sathyaprakash ---
% \textit{Dark Sirens to Resolve the Hubble–Lemaître Tension}\\
% {\small \textit{Astrophys.~J.~Lett.}~905~(2020)~2,~L28
% \hfill \href{https://arxiv.org/abs/2007.02883}{arXiv:2007.02883}}

% \vspace{10pt}
% {\small\color{muted}\textit{Full list of publications (incl.\ LIGO, Cosmic Explorer, 
% and Einstein Telescope collaboration papers):}\\
% \href{https://arxiv.org/search/?searchtype=author&query=borhanian}{\texttt{arxiv.org/search/?searchtype=author\&query=borhanian}}}

%──────────────────────────────────────────────────────────────────────────────
\end{document}
